\begin{figure*}[t]
     \centering
     \includegraphics[trim=0 0 0 0,clip, width=1\linewidth]{Figures/main.pdf}
     \setlength{\abovecaptionskip}{-5pt}
     \setlength{\belowcaptionskip}{-10pt}
     \caption{\textbf{Overview of \OURMODEL.}
    Given multi-view context images, the Geometry Backbone predicts camera parameters, multi-scale Gaussian parameter maps, and densification score maps. During training, camera and Gaussian parameters are jointly optimized with the camera loss $\mathcal{L}^{\text{camera}}$, rendering loss $\mathcal{L}^{\text{render}}$, and scene-scale regularization $\mathcal{L}^{\text{scene}}$. The predicted densification score maps are trained by score loss $\mathcal{L}^{\text{score}}$ derived from the backpropagation of rendering loss. The final representation $\mathcal{G}_{\tau}$, constructed using a randomly sampled threshold $\tau$, is also supervised by the rendering loss $\mathcal{L}^{\text{render}}$. 
    At inference time, the predicted Gaussian parameter maps and densification scores are used to generate compact, high-fidelity 3D Gaussian representations $\mathcal{G}_{\tau_{\bar{N}_{\mathcal{G}}}}$ tailored to user-specified Gaussian budgets $\bar{N}_{\mathcal{G}}$. This adaptation is efficient and does not require retraining.
     }
     \label{fig:main}
\end{figure*}